\clearpage
\section{선별 문제 풀이 I}
\label{sec:abel-ruffini-solutions-a}

\subsection*{문제 \thechapter.1: $A_5$의 켤레류}

\begin{solution}
$A_5$의 비항등원은 순환형 $3+1+1$, $2+2+1$, $5$ 가운데 하나이다.

삼순환은
\[
  \binom53\cdot2=20
\]
개이다. 이중전치는 고정점을 $5$가지로 고르고 나머지 네 점을 두 쌍으로 나누는 $3$가지 방법이 있으므로 $15$개이다. 오순환은 $(5-1)!=24$개이다.

$\alpha=(1\,2\,3)$의 $S_5$-중심화군은 $\langle\alpha\rangle\times\langle(4\,5)\rangle$이고, 그 짝순열 부분의 위수는 $3$이다. 따라서
\[
  |\operatorname{Cl}_{A_5}(\alpha)|=60/3=20.
\]

$\beta=(1\,2)(3\,4)$의 $S_5$-중심화군 위수는 $8$이고, 그 짝순열 부분은 $V_4$이므로
\[
  |\operatorname{Cl}_{A_5}(\beta)|=60/4=15.
\]

$\gamma=(1\,2\,3\,4\,5)$의 중심화군은 $\langle\gamma\rangle$이고 위수는 $5$이다. 따라서 각 오순환 켤레류의 크기는 $60/5=12$이다. 오순환은 $24$개이므로 두 류로 갈라진다. 결론적으로 켤레류 크기는
\[
  \boxed{1,20,15,12,12}
\]
이다.
\end{solution}

\subsection*{문제 \thechapter.3: $A_5$의 단순성}

\begin{solution}
$N\triangleleft A_5$이면 $N$은 항등원 켤레류와 다른 켤레류들 가운데 일부의 합집합이다. 따라서 가능한 위수는
\[
  1+\text{일부의 }20,15,12,12
\]
이다. 모든 합을 적으면
\[
  1,13,16,21,25,28,33,36,40,45,48,60.
\]
라그랑주 정리에 의해 $|N|$은 $60$의 약수여야 한다. 위 목록에서 $60$의 약수는 $1$과 $60$뿐이다. 따라서 $N=1$ 또는 $N=A_5$이고, $A_5$는 단순군이다.
\end{solution}

\subsection*{문제 \thechapter.5: $A_n$의 단순성}

\begin{solution}
$1\ne N\triangleleft A_n$에서 $1\ne\sigma\in N$을 잡는다. 정상성 때문에 임의의 $\tau\in A_n$에 대해 $[\tau,\sigma]\in N$이다.

$\sigma$가 길이 $4$ 이상인 순환 $(a\,b\,c\,d\,\dots)$를 포함하면
\[
  [(a\,b\,c),\sigma]=(a\,b\,d)
\]
이므로 $N$은 삼순환을 포함한다.

$\sigma$가 삼순환 $(a\,b\,c)$을 포함한다고 하자. 이것이 $\sigma$ 전체가 아니면 다른 움직이는 점 $d$와 $e=\sigma(d)$를 택할 수 있고
\[
  [(a\,b\,d),\sigma]=(a\,b\,e\,c\,d)
\]
인 오순환을 얻는다. 앞 경우를 이 오순환에 적용하면 삼순환을 얻는다.

마지막으로 $\sigma$가 전치들의 곱이라고 하자. 고정점 $e$가 있으면 전치 $(a\,b)$에 대해
\[
  [(a\,b\,e),\sigma]=(a\,e\,b)
\]
이다. 고정점이 없으면 전치의 개수는 짝수이고 적어도 네 개이다. 세 전치를 $(a\,b)(c\,d)(e\,f)$라 쓰면
\[
  [(a\,c\,e),\sigma]=(a\,c\,e)(b\,f\,d).
\]
이는 두 삼순환의 곱이므로 두 번째 경우를 적용해 결국 삼순환을 얻는다.

$n\ge5$에서 모든 삼순환은 $A_n$ 안에서 서로 켤레이고, 삼순환들이 $A_n$을 생성한다. 따라서 $N=A_n$이다.
\end{solution}

\subsection*{문제 \thechapter.7: $S_n$과 $A_n$의 도함열}

\begin{solution}
부호준동형의 핵이 $A_n$이므로 $S_n/A_n$은 아벨군이고
\[
  [S_n,S_n]\subseteq A_n.
\]
모든 삼순환은 두 전치의 교환자이며 삼순환들이 $A_n$을 생성하므로 반대 포함도 성립한다.
\[
  [S_n,S_n]=A_n.
\]

$n\ge5$에서 서로 다른 $a,b,c,d,e$에 대해
\[
  [(a\,b\,d),(a\,c\,e)]=(a\,b\,c).
\]
따라서 모든 삼순환이 $[A_n,A_n]$에 들어가고
\[
  [A_n,A_n]=A_n.
\]
그러므로
\[
  S_n^{(0)}=S_n,
  \qquad
  S_n^{(r)}=A_n\quad(r\ge1).
\]
도함열은 $1$에 도달하지 않으므로 $S_n$과 $A_n$은 비가해이다.

또한 $A_n$의 단순성을 사용하면 $S_n$의 정상부분군은 $1,A_n,S_n$뿐이다. 실제로 $N\triangleleft S_n$이면 $N\cap A_n$은 $1$ 또는 $A_n$이고, 첫 경우에는 부호준동형이 $N$에 단사이므로 $|N|\le2$이다. 비자명한 위수 $2$의 정상부분군은 중심에 들어가지만 $Z(S_n)=1$이므로 불가능하다.
\end{solution}

\subsection*{문제 \thechapter.9: $A_5$의 Sylow 부분군}

\begin{solution}
위수 $5$인 부분군의 비항등원은 오순환 네 개이다. 오순환은 모두 $24$개이고 서로 다른 위수 $5$ 부분군은 항등원 밖에서 만나지 않으므로
\[
  n_5=24/4=6.
\]

위수 $3$인 부분군의 비항등원은 서로 역인 삼순환 두 개이다. 삼순환은 $20$개이므로
\[
  n_3=20/2=10.
\]

$|A_5|=60$에서 $2$-Sylow 부분군의 위수는 $4$이다. $A_5$에는 위수 $4$ 원소가 없으므로 각 $2$-Sylow 부분군은 $V_4$이다. 각 $V_4$는 비항등 이중전치 세 개를 포함하고, 한 이중전치와 항등원은 그 이중전치와 함께 같은 네 점 위의 나머지 두 이중전치를 유일하게 결정한다. 따라서
\[
  n_2=15/3=5.
\]
Sylow 합동조건과도 일치한다.
\end{solution}

\subsection*{문제 \thechapter.11: 대칭 유리함수의 고정체}

\begin{solution}
$K=L^{S_n}$라 하자. 아르틴의 고정체 정리에서
\[
  [L:K]=|S_n|=n!.
\]
기본대칭식들은 $S_n$에 고정되므로
\[
  k(s_1,\dots,s_n)\subseteq K.
\]

다항식
\[
  f(X)=\prod_{i=1}^n(X-T_i)
  =X^n-s_1X^{n-1}+\cdots+(-1)^ns_n
\]
의 분해체가 $L$이다. 분해체의 차수는 많아야 $n!$이므로
\[
  [L:k(s_1,\dots,s_n)]\le n!.
\]
한편 중간체 차수 공식에서
\[
  [L:k(s_1,\dots,s_n)]
  =[L:K][K:k(s_1,\dots,s_n)]
  \ge n!.
\]
따라서 두 부등식이 모두 등호이고
\[
  \boxed{L^{S_n}=k(s_1,\dots,s_n)}.
\]
\end{solution}

\subsection*{문제 \thechapter.13: 일반다항식의 갈루아군}

\begin{solution}
$s_1,\dots,s_n$은 $k$ 위에서 대수적으로 독립이므로
\[
  \varphi:k(C_1,\dots,C_n)
  \longrightarrow k(s_1,\dots,s_n),
  \qquad
  C_j\longmapsto(-1)^js_j
\]
는 체동형이다. 이 동형 아래에서
\[
  X^n+C_1X^{n-1}+\cdots+C_n
\]
은
\[
  X^n-s_1X^{n-1}+s_2X^{n-2}-\cdots+(-1)^ns_n
  =\prod_{i=1}^n(X-T_i)
\]
로 옮겨진다.

후자의 분해체는 $k(T_1,\dots,T_n)$이고 고정체 계산에서 갈루아군은 변수들을 임의로 순열하는 $S_n$이다. 따라서
\[
  \boxed{\Gal(P_n/k(C_1,\dots,C_n))\cong S_n}.
\]
추이성에서 기약성이, 서로 다른 독립변수라는 사실에서 분리성이 따른다.
\end{solution}

\subsection*{문제 \thechapter.15: 일반 삼차식의 판별식 층}

\begin{solution}
일반 삼차식의 갈루아군은 $S_3$이다. 반데르몬드 곱
\[
  \delta=(T_1-T_2)(T_1-T_3)(T_2-T_3)
\]
는 짝순열에 고정되고 홀순열에 의해 부호가 바뀐다. 따라서
\[
  K_3(\delta)\subseteq L^{A_3}.
\]
$\charac k\ne2$이므로 $\delta\notin K_3$이고 왼쪽 확장의 차수는 $2$이다. 오른쪽도 $[S_3:A_3]=2$이므로
\[
  L^{A_3}=K_3(\delta)
  =K_3(\sqrt{\Disc(P_3)}).
\]

정규열
\[
  S_3\trianglerighteq A_3\trianglerighteq1
\]
의 인자군은 $C_2,C_3$이다. 첫 단계는 판별식의 제곱근을 첨가하는 이차층이고, 원시 세제곱근을 준비한 뒤 둘째 단계는 Lagrange 분해식의 세제곱근을 첨가하는 순환 삼차층이다. 이것이 Cardano 공식의 체론적 구조이다.
\end{solution}
